\documentclass[12pt]{article}
\usepackage[margin=0.72in]{geometry}
\usepackage{lmodern}
\usepackage{microtype}
\usepackage{multicol}
\usepackage{fancyhdr}
\usepackage{titlesec}
\usepackage{enumitem}
\usepackage{xcolor}
\usepackage[hidelinks]{hyperref}

\setlength{\columnsep}{0.30in}
\setlength{\parindent}{1.05em}
\setlength{\parskip}{0.24em}
\setlength{\headheight}{15pt}
\titleformat{\section}{\large\bfseries\scshape}{}{0pt}{}
\titleformat{\subsection}{\normalsize\bfseries}{}{0pt}{}
\pagestyle{fancy}
\fancyhf{}
\fancyhead[L]{Douglass and Literacy}
\fancyhead[R]{Reading Handout}
\fancyfoot[C]{\thepage}
\renewcommand{\headrulewidth}{0.4pt}

\newcommand{\focusbox}[1]{%
\vspace{0.4em}\noindent\colorbox{gray!12}{\begin{minipage}{0.96\linewidth}#1\end{minipage}}\vspace{0.5em}}

\begin{document}
\begin{center}
{\LARGE\bfseries Do We Love Shadows?}\par\vspace{0.25em}
{\large\scshape Frederick Douglass and Literacy as Liberation}\par\vspace{0.45em}
\rule{0.82\textwidth}{0.4pt}\par\vspace{0.7em}
\end{center}

\section*{Vocabulary Preview}
\begin{multicols}{2}
\begin{itemize}[leftmargin=*, itemsep=0.18em]
\item \textbf{Literacy}: The ability to read and write.
\item \textbf{Enslaved}: Held as property and denied freedom.
\item \textbf{Enslaver}: A person who claims ownership and power over an enslaved person.
\item \textbf{Liberation}: Being set free.
\item \textbf{Ignorance}: Lack of knowledge or understanding.
\item \textbf{Tyranny}: Cruel or unjust power over others.
\item \textbf{Dehumanizing}: Treating a person as less than fully human.
\item \textbf{Stratagem}: A plan or trick used to overcome difficulty.
\item \textbf{Abolition}: The movement to end slavery.
\item \textbf{Consciousness}: Awareness of oneself and one's condition.
\item \textbf{Discontented}: Unsatisfied; unwilling to accept one's present condition.
\item \textbf{Columbian Orator}: A schoolbook Douglass used that contained speeches about liberty and human rights.
\end{itemize}
\end{multicols}

\section*{Introduction}
Frederick Douglass was born into slavery in Maryland around 1818 and escaped to freedom in 1838. He became one of the most powerful writers and speakers against slavery in American history. In his \textit{Narrative}, Douglass explains how learning to read changed his understanding of himself, slavery, and freedom.

Douglass helps us see that shadows are not always accidental. Plato shows prisoners trapped in a cave; Bacon shows false images inside the mind; Augustine shows a heart divided by old loves. Douglass adds another truth: unjust power may deliberately keep people in darkness because knowledge can awaken a desire for freedom. As you read, ask: \textbf{Why would someone in power want others kept in darkness?}

\focusbox{\textbf{What you need to know before this passage:} As a child, Douglass is sent from a plantation to Baltimore to live with Hugh and Sophia Auld. At first, Sophia Auld is unusually kind to him because she has not previously been trained by slavery's habits of power. She begins teaching him the alphabet. Her husband stops her, and his explanation reveals to Douglass the connection between literacy and freedom.}

\focusbox{\textbf{Source note:} Adapted and modernized from Frederick Douglass, \textit{Narrative of the Life of Frederick Douglass, an American Slave}, chapters VI--VII. This handout uses the focused literacy episode: Sophia Auld's first lessons, Hugh Auld's command to stop them, Douglass's secret education, and the painful awakening that followed.}

\begin{multicols}{2}
\section*{Reading}

\subsection*{1. A Mistress Not Yet Trained by Slavery}
My new mistress proved to be all she appeared when I first met her at the door: a woman with a kind heart and fine feelings. She had never before had an enslaved person under her control. Before her marriage, she had depended on her own work for a living. She was a weaver, and by steady work she had been partly preserved from the blighting and dehumanizing effects of slavery.

I was utterly astonished at her goodness. I scarcely knew how to behave toward her. She was entirely unlike any other white woman I had ever seen. I could not approach her as I had been taught to approach other white ladies. The crouching servility usually expected from an enslaved person did not please her; it disturbed her. She did not think it impudent for an enslaved person to look her in the face. The meanest enslaved person could put her at ease, and none left her presence without feeling better.

Her face was made of heavenly smiles, and her voice was gentle music. But, sadly, this kind heart did not remain unchanged. Slavery soon proved able to take away these excellent qualities. The poisonous power was already in her hands, and it soon began its terrible work. Her cheerful eye, under slavery's influence, became red with rage. Her sweet voice changed into harsh discord. Her angelic face gave place to the face of a demon.

\subsection*{2. The Forbidden Lesson}
Very soon after I came to live with Mr. and Mrs. Auld, she kindly began to teach me the A, B, C. After I learned the alphabet, she helped me learn to spell words of three or four letters.

At this point, Mr. Auld discovered what was happening. He immediately forbade Mrs. Auld to teach me any further. He told her that it was unlawful and unsafe to teach an enslaved person to read. To use the substance of his words, he said that an enslaved person should know nothing except to obey his master and do as he is told. Learning, he said, would spoil the best enslaved person in the world.

He said that if I learned to read, there would be no keeping me. It would forever unfit me to be a slave. I would at once become unmanageable and of no value to my master. As for me, he said, reading could do me no good, but a great deal of harm. It would make me discontented and unhappy.

These words sank deep into my heart. They awakened feelings that slept before and called into existence an entirely new train of thought. It was a new and special revelation, explaining dark and mysterious things with which my young understanding had struggled but struggled in vain.

I now understood what had been to me a most puzzling difficulty: the white man's power to enslave the black man. From that moment, I understood the pathway from slavery to freedom.

Though I understood the difficulty of learning without a teacher, I set out with high hope and a fixed purpose, whatever the cost. In learning to read, I owed almost as much to my master's bitter opposition as to my mistress's first kindness. I admit the benefit of both.

\subsection*{3. The Change in Sophia Auld}
My mistress had given me the first inch by teaching me the alphabet. No precaution could prevent me from taking the whole ell.

But after her husband's command, she changed. At first she had lacked the depravity necessary to shut me up in mental darkness. She had treated me as she thought one human being ought to treat another. Slavery could not bear such kindness. Under its influence, the tender heart became stone. The lamb-like nature gave way to tiger-like fierceness.

Nothing seemed to make her angrier than seeing me with a newspaper. She seemed to think that if I had a newspaper, I must be learning to read. From that time, I was watched closely. If I spent much time in a separate room, I was suspected of having a book and called to explain myself.

All this, however, was too late. The first step had already been taken.

\subsection*{4. Bread for Knowledge}
The plan by which I was most successful was to make friends of all the little white boys I met in the street. As many as I could, I turned into teachers. With their kind help, at different times and in different places, I finally learned to read.

When I was sent on errands, I always took my book with me. By doing one part of the errand quickly, I found time to get a lesson before returning. I also carried bread with me. There was usually enough bread in the house, and I was welcome to it. In this way, I was better off than many poor white children in our neighborhood.

I gave bread to the hungry boys, and they gave me that more valuable bread of knowledge.

I am strongly tempted to name some of those boys, but prudence forbids it. It would embarrass them, for it was almost an unforgivable offense to teach an enslaved person to read in that Christian country. I will not name them, but they have my love and gratitude.

I often told them that I wished I could be as free as they would be when they became men. They would tell me I would be free someday. I did not believe them, but still I hoped.

\subsection*{5. A Book Opens the Argument}
At about the age of twelve, I got hold of a book called \textit{The Columbian Orator}. Every chance I had, I read it. In that book, I found a dialogue between a master and an enslaved person. The enslaved person had run away three times. In the dialogue, the master argued for slavery, but the enslaved person answered him. He exposed the arguments for slavery and argued for human rights. In the end, the master was convinced and freed him.

The book also contained one of Sheridan's speeches on behalf of Catholic emancipation. These writings gave speech to interesting thoughts in my own soul. They gave me powerful arguments against slavery, and at the same time they made me hate my enslavers even more.

What I gained from reading was an understanding of the evil system under which I suffered. I could now see that slavery was not merely my private hardship. It was a system of robbery, violence, and falsehood. Reading gave me words for the deep feelings I already had.

\subsection*{6. The Pain of Awakening}
Yet the more I read, the more I was led to hate and detest those who enslaved me. I could regard them only as a band of successful robbers who had left their homes, gone to Africa, stolen us from our homes, and reduced us to slavery in a strange land.

I hated them as the meanest and most wicked of men. As I read and thought, I became miserable. I sometimes felt that learning to read had been a curse rather than a blessing. It had given me a view of my wretched condition, but no remedy.

It opened my eyes to the horrible pit, but it showed no ladder by which to climb out. In moments of agony, I envied my fellow enslaved people for their ignorance. I wished myself a beast. Anything, no matter what, seemed better than thinking about my condition.

The silver trumpet of freedom had awakened my soul to eternal wakefulness. Freedom had now appeared, never to disappear again. It was heard in every sound and seen in every thing. It was always present to torment me with a sense of my miserable condition.

I saw nothing without seeing freedom. I heard nothing without hearing freedom. It looked from every star. It smiled in every calm, breathed in every wind, and moved in every storm.

\subsection*{7. Learning to Write}
Meanwhile, I learned to write. In the shipyard, I noticed that the carpenters marked pieces of timber with letters to show where they belonged. The part for the larboard side was marked L. The starboard side was marked S. The larboard forward piece was marked L. F.; the larboard aft piece, L. A.; the starboard forward piece, S. F.; and the starboard aft piece, S. A. I soon learned the names of these letters and what they meant.

When I met any boy who could write, I would tell him I could write as well as he could. He would say I could not. I would then make the letters I already knew and ask him to beat that. In this way, I learned new letters. By copying from books and writing between the lines of old copybooks, I slowly taught myself to write.

Thus, after a long and tedious effort for years, I finally succeeded in learning how to write.

\section*{Reading Focus}
\begin{enumerate}[leftmargin=*, itemsep=0.2em]
\item Underline Mr. Auld's explanation for why Douglass must not learn to read.
\item Circle one sentence showing that knowledge can be painful.
\item Put a star beside one sentence showing that literacy leads Douglass toward freedom.
\item Write ``power'' in the margin beside one place where someone tries to control knowledge.
\item Box one sentence that best connects Douglass back to Plato's cave.
\end{enumerate}

\section*{Power and Knowledge Map}
Create a two-sided map in your notebook.

\begin{itemize}[leftmargin=*, itemsep=0.2em]
\item \textbf{Side 1: How slavery uses darkness} --- List three ways ignorance helps preserve unjust power.
\item \textbf{Side 2: How literacy brings light} --- List three ways reading or writing helps Douglass move toward freedom.
\end{itemize}

\section*{Connection to Plato}
Create a comparison with two parts: \textbf{Plato's Cave} and \textbf{Douglass's Narrative}. For each one, answer: Who is trapped? What are the shadows? Why is truth painful? What does freedom require?

\section*{Claim Question}
Answer in one clear paragraph: \textbf{Why would unjust power try to keep people ignorant?} State a clear claim, use one example from Douglass, explain how ignorance protects power, and connect Douglass to Plato, Bacon, or Augustine.

\section*{Handout Summary}
Finish by writing a short summary of this handout. In three to five sentences, explain how Douglass learns to read and write, why that knowledge is dangerous to slavery, and how this reading changes the unit's idea of shadows.

\end{multicols}
\end{document}
